\section{Related Work}
\label{sec:related}

\paragraph{Harness optimization and learned editors.}
DSPy compiles declarative pipelines of LM calls \citep{khattab2024dspy}, and AFlow searches over
code-represented workflows \citep{zhang2025aflow}. Meta-Harness and AutoHarness revise or synthesize
harness code from execution feedback \citep{lee2026metaharness,lou2026autoharness}. Harness-R1 trains an
editor with SFT and GRPO against a frozen target agent \citep{shao2026harnessr1}, and JIT-Agent learns
task-conditioned harness generation, repair and evolution with rewards that also account for latency and
cost \citep{zhang2026jitagent}. Harness learning \citep{zhang2026harnesslearning} casts revision as
meta-learning over executable programs and is the starting point of this paper. These works focus on
\emph{producing} better revisions. We focus on \emph{accepting} them, and our policies can be placed
after any of these proposers.

\paragraph{Evaluation of harness evolution.}
Several 2026 studies question how harness evolution is evaluated. \citet{wang2026rethinking} examine
benchmark overfitting under matched feedback and inference budgets. \citet{lin2026updating} separate
producing an update from benefiting from it. HarnessCompass constrains edits to improve transfer
\citep{zhang2026harnesscompass}, and Evo-Bench evaluates evolution across held-out task suites
\citep{huang2026evobench}. HarnessDev evaluates harness creation and evolution with execution-token
cost and reports unstable gains that only partly transfer \citep{wu2026harnessdev}. EvoHarnessBench
documents \emph{harness-induced forgetting}, where expanding a harness lowers performance on tasks that
were already solved \citep{ke2026evoharnessbench}. Our matched-budget protocol and regression
measurements follow these concerns. We do not claim either as new.

\paragraph{Gated self-evolution.}
HarnessEvolve is closest to our work \citep{harnessevolve2026}. It combines a quality gate, an
isolated agent that checks for data leakage and prompt bloat, with a performance gate that requires
improvement on the current batch without degradation on recent ones. \citet{cai2026reliable} argue on theoretical grounds that
improving previously failed tasks does not guarantee that earlier successes are preserved. They propose a
condition for adopting an evolved harness and a validation rule that enforces it. Our regression gate is
an empirical, finite-sample instance of the same concern. RAVEL differs from both in three ways. It replaces point-estimate
gates with \emph{finite-sample statistical} ones: a Wilson bound on regressions, an upper bound on cost,
and an exact paired test for improvement. It adds runtime security probes and fault-aware comparison.
And it is evaluated by holding the proposer fixed and varying only the acceptance rule. We regard
these as complementary designs rather than competing systems, and we do not reproduce HarnessEvolve's
results.

\paragraph{Self-improving agents and safety.}
Reflexion and Self-Refine improve outputs from feedback \citep{shinn2023reflexion,madaan2023selfrefine}.
STOP and the Darwin G\"odel Machine let agents modify their own code
\citep{zelikman2024stop,zhang2025dgm}, and the DGM authors note the need for sandboxing and human
oversight. Indirect prompt injection, where instructions are hidden in retrieved content, is a
well-documented threat to tool-using agents \citep{greshake2023injection}. Our adversarial-proposer
condition goes further and assumes the \emph{proposer itself} may emit harmful code, a threat model
appropriate whenever revisions come from a model whose outputs are not fully trusted.

\paragraph{Statistical acceptance testing.}
We use the Wilson score interval \citep{wilson1927} for regression rates and the exact McNemar (sign)
test \citep{mcnemar1947} for paired improvements. Both are standard. What we contribute is applying
them as acceptance gates in a harness-evolution loop and measuring what they cost.
